\documentclass[12pt,a4paper]{article}

\usepackage[T2A]{fontenc}
\usepackage[utf8]{inputenc}
\usepackage[russian]{babel}
\usepackage{amsmath,amssymb,amsthm}
\usepackage{mathtools}
\usepackage{enumitem}
\usepackage{listings}
\usepackage{xcolor}
\usepackage{array}
\usepackage{hyperref}
\usepackage{stmaryrd}

\lstdefinestyle{code}{
  basicstyle=\ttfamily\small,
  keywordstyle=\color{blue},
  commentstyle=\color{gray},
  stringstyle=\color{red!60!black},
  columns=fullflexible,
  keepspaces=true,
  frame=single,
  breaklines=true,
  showstringspaces=false
}

\lstset{style=code}

\newtheorem{definition}{Определение}
\newtheorem{theorem}{Теорема}
\newtheorem{lemma}{Лемма}
\newtheorem{example}{Пример}
\newtheorem{remark}{Замечание}

\title{Билет №33 \\ \small Языки программирования со статической и динамической типизацией. Сильная и слабая типизация. Языки сценариев. (СпБГУ)}
\author{Сериков Владислав Александрович}
\date{13 мая 2026}

\begin{document}

\maketitle
\tableofcontents
\newpage

\section{Языки программирования со статической и динамической типизацией.
Сильная и слабая типизация. Языки сценариев}

\subsection{Базовые понятия}

\textbf{Тип} --- это множество значений вместе с набором допустимых операций над ними.
Например, тип \texttt{int} обычно задаёт множество целых чисел некоторого диапазона
и операции сложения, вычитания, умножения, сравнения и т.д.

\textbf{Система типов} --- это совокупность правил языка программирования, которые
назначают типы выражениям, переменным, функциям и другим программным конструкциям,
а также определяют, какие операции над ними допустимы.

Основная задача системы типов --- предотвращать часть некорректных программ ещё до
или во время выполнения. Типичные ошибки, связанные с типами:

\begin{itemize}
    \item сложение числа и строки без явно определённого смысла этой операции;
    \item вызов функции с аргументом неподходящего типа;
    \item обращение к полю, которого нет у объекта;
    \item использование значения как функции, если оно функцией не является;
    \item обращение с указателем или ссылкой как со значением другого типа.
\end{itemize}

Например:

\begin{verbatim}
x = "hello"
y = x + 1
\end{verbatim}

В одних языках такая программа будет отвергнута как ошибка типов, в других язык
попытается выполнить неявное преобразование, например превратить число \texttt{1}
в строку \texttt{"1"}.

\subsection{Статическая и динамическая типизация}

\subsubsection{Статическая типизация}

\textbf{Статическая типизация} --- это способ организации системы типов, при котором
типы выражений проверяются до выполнения программы, обычно на этапе компиляции.

Иначе говоря, если язык статически типизирован, то значительная часть ошибок типов
обнаруживается до запуска программы.

Примеры языков со статической типизацией:

\begin{itemize}
    \item C;
    \item C++;
    \item Java;
    \item C\#;
    \item Rust;
    \item Go;
    \item Haskell;
    \item OCaml;
    \item Scala;
    \item Kotlin;
    \item TypeScript.
\end{itemize}

Пример на Java:

\begin{verbatim}
int x = 10;
x = "hello"; // ошибка компиляции
\end{verbatim}

Переменная \texttt{x} объявлена как \texttt{int}, поэтому ей нельзя присвоить строку.

\subsubsection{Динамическая типизация}

\textbf{Динамическая типизация} --- это способ организации системы типов, при котором
типы значений проверяются во время выполнения программы.

В динамически типизированных языках переменная обычно не имеет фиксированного типа,
но каждое значение имеет тип. Одна и та же переменная может в разные моменты времени
ссылаться на значения разных типов.

Примеры языков с динамической типизацией:

\begin{itemize}
    \item Python;
    \item JavaScript;
    \item Ruby;
    \item PHP;
    \item Lua;
    \item Perl;
    \item Tcl;
    \item Lisp;
    \item Scheme;
    \item Smalltalk.
\end{itemize}

Пример на Python:

\begin{verbatim}
x = 10
x = "hello"  # допустимо
\end{verbatim}

Здесь имя \texttt{x} сначала связано с целым числом, а затем со строкой.

Однако операция с неподходящими типами вызовет ошибку во время выполнения:

\begin{verbatim}
x = "hello"
y = x + 1  # TypeError во время выполнения
\end{verbatim}

\subsubsection{Сравнение статической и динамической типизации}

\begin{center}
\begin{tabular}{|p{0.28\textwidth}|p{0.32\textwidth}|p{0.32\textwidth}|}
\hline
\textbf{Критерий} & \textbf{Статическая типизация} & \textbf{Динамическая типизация} \\
\hline
Момент проверки типов &
До выполнения программы &
Во время выполнения программы \\
\hline
Тип переменной &
Обычно фиксируется явно или выводится компилятором &
Переменная может ссылаться на значения разных типов \\
\hline
Обнаружение ошибок &
Многие ошибки находятся до запуска &
Часть ошибок проявляется только при выполнении \\
\hline
Гибкость &
Ниже, но выше надёжность интерфейсов &
Выше, удобно для быстрой разработки \\
\hline
Производительность &
Часто выше за счёт знания типов заранее &
Часто ниже из-за проверок во время выполнения, хотя JIT-компиляция может снижать разницу \\
\hline
Примеры &
C, C++, Java, Rust, Haskell &
Python, JavaScript, Ruby, Lua \\
\hline
\end{tabular}
\end{center}

\subsection{Явная и неявная типизация}

Статическую и динамическую типизацию не следует смешивать с явной и неявной
типизацией.

\textbf{Явная типизация} означает, что программист явно указывает типы в программе.

Пример на Java:

\begin{verbatim}
int x = 10;
String s = "hello";
\end{verbatim}

\textbf{Неявная типизация} означает, что типы выводятся автоматически из контекста.

Пример на C++:

\begin{verbatim}
auto x = 10;        // x имеет тип int
auto s = "hello";   // s имеет тип const char*
\end{verbatim}

Пример на Haskell:

\begin{verbatim}
square x = x * x
\end{verbatim}

Компилятор Haskell способен вывести тип функции без явного указания программистом.

Важно:

\begin{itemize}
    \item язык может быть статически типизированным и иметь вывод типов;
    \item язык может быть динамически типизированным и при этом позволять аннотации типов;
    \item явность типов и момент проверки типов --- разные свойства.
\end{itemize}

Например, Haskell является статически типизированным языком с мощным выводом типов,
а TypeScript является статически типизированным расширением JavaScript, где многие
типы также могут выводиться автоматически.

\subsection{Сильная и слабая типизация}

\subsubsection{Проблема терминологии}

Термины \textbf{сильная типизация} и \textbf{слабая типизация} менее формальны, чем
термины \textbf{статическая} и \textbf{динамическая типизация}. В литературе и практике
они могут использоваться по-разному.

В общем виде:

\begin{itemize}
    \item \textbf{сильно типизированный язык} строго ограничивает операции над
    значениями разных типов и не выполняет опасных неявных преобразований;
    \item \textbf{слабо типизированный язык} допускает более свободные неявные
    преобразования типов или позволяет интерпретировать данные одного типа как данные
    другого типа.
\end{itemize}

Степень ``силы'' типизации является скорее шкалой, чем бинарным свойством.

\subsubsection{Сильная типизация}

\textbf{Сильная типизация} означает, что язык не позволяет выполнять операции над
значениями неподходящих типов без явного преобразования или без явно заданного
смысла операции.

Пример на Python:

\begin{verbatim}
"10" + 5
\end{verbatim}

Результат:

\begin{verbatim}
TypeError
\end{verbatim}

Python является динамически типизированным, но обычно считается сильно
типизированным языком, поскольку он не выполняет произвольные неявные преобразования
между строками и числами в арифметических операциях.

Корректный вариант:

\begin{verbatim}
int("10") + 5      # 15
"10" + str(5)     # "105"
\end{verbatim}

Здесь преобразование типов выполняется явно.

\subsubsection{Слабая типизация}

\textbf{Слабая типизация} означает, что язык допускает неявные преобразования,
которые могут быть неожиданными для программиста.

Пример на JavaScript:

\begin{verbatim}
"10" + 5    // "105"
"10" - 5    // 5
\end{verbatim}

В первом случае число \texttt{5} преобразуется в строку, поскольку оператор
\texttt{+} используется также для конкатенации строк. Во втором случае строка
\texttt{"10"} преобразуется в число, поскольку оператор \texttt{-} не имеет смысла
для строковой конкатенации.

Ещё один пример:

\begin{verbatim}
0 == false     // true
0 === false    // false
\end{verbatim}

Оператор \texttt{==} выполняет неявные преобразования типов, а оператор \texttt{===}
сравнивает без таких преобразований.

\subsubsection{Слабая типизация в C}

Язык C обычно считается статически типизированным, но в ряде аспектов слабо
типизированным или допускающим низкоуровневые операции, ослабляющие контроль типов.

Пример:

\begin{verbatim}
int x = 65;
char c = (char)x;
\end{verbatim}

Здесь целое число явно приводится к символу.

Более опасный пример связан с указателями:

\begin{verbatim}
int x = 42;
void *p = &x;
double *q = (double *)p;
\end{verbatim}

Язык позволяет преобразовать указатель одного типа к указателю другого типа, хотя
использование такого значения может привести к неопределённому поведению.

\subsubsection{Сравнение свойств}

Статическая/динамическая типизация и сильная/слабая типизация --- независимые
характеристики.

\begin{center}
\begin{tabular}{|p{0.27\textwidth}|p{0.32\textwidth}|p{0.32\textwidth}|}
\hline
 & \textbf{Сильная типизация} & \textbf{Слабая типизация} \\
\hline
\textbf{Статическая типизация} &
Java, Haskell, Rust, OCaml &
C, C++ в некоторых аспектах \\
\hline
\textbf{Динамическая типизация} &
Python, Ruby, Scheme &
JavaScript, PHP в некоторых аспектах \\
\hline
\end{tabular}
\end{center}

Эта таблица является приближённой: языки могут иметь как строгие, так и слабые
механизмы в разных частях системы типов.

\subsection{Проверка типов}

\subsubsection{Задача проверки типов}

\textbf{Проверка типов} --- это процесс установления того, удовлетворяет ли программа
правилам системы типов языка.

Для статически типизированного языка проверка типов выполняется до запуска программы.
Для динамически типизированного языка проверки выполняются во время исполнения.

Типовая проверка обычно решает следующие задачи:

\begin{itemize}
    \item сопоставить каждому выражению тип;
    \item проверить корректность операций;
    \item проверить корректность вызовов функций;
    \item проверить корректность присваиваний;
    \item проверить корректность возвращаемых значений;
    \item проверить обращения к полям и методам объектов.
\end{itemize}

\subsubsection{Минимальный алгоритм статической проверки типов}

Рассмотрим простой язык с выражениями:

\[
e ::= n \mid \texttt{true} \mid \texttt{false} \mid x
    \mid e_1 + e_2
    \mid e_1 < e_2
    \mid \texttt{if } e_1 \texttt{ then } e_2 \texttt{ else } e_3.
\]

Типы:

\[
\tau ::= \texttt{Int} \mid \texttt{Bool}.
\]

Пусть имеется окружение типов:

\[
\Gamma : \text{Var} \to \{\texttt{Int}, \texttt{Bool}\},
\]

где \(\Gamma(x)\) --- тип переменной \(x\).

Алгоритм проверки типов выражения \(e\):

\begin{enumerate}
    \item Если \(e\) является целочисленной константой \(n\), вернуть тип
    \(\texttt{Int}\).
    \item Если \(e\) равно \texttt{true} или \texttt{false}, вернуть тип
    \(\texttt{Bool}\).
    \item Если \(e\) является переменной \(x\), найти \(x\) в окружении \(\Gamma\).
    Если переменная не найдена, сообщить об ошибке.
    \item Если \(e = e_1 + e_2\), рекурсивно проверить типы \(e_1\) и \(e_2\).
    Оба должны иметь тип \(\texttt{Int}\). Тогда результат имеет тип \(\texttt{Int}\).
    \item Если \(e = e_1 < e_2\), рекурсивно проверить типы \(e_1\) и \(e_2\).
    Оба должны иметь тип \(\texttt{Int}\). Тогда результат имеет тип \(\texttt{Bool}\).
    \item Если
    \[
    e = \texttt{if } e_1 \texttt{ then } e_2 \texttt{ else } e_3,
    \]
    то:
    \begin{enumerate}
        \item проверить, что \(e_1\) имеет тип \(\texttt{Bool}\);
        \item проверить типы \(e_2\) и \(e_3\);
        \item убедиться, что типы ветвей совпадают;
        \item вернуть общий тип ветвей.
    \end{enumerate}
\end{enumerate}

\subsubsection{Иллюстрация работы алгоритма}

Рассмотрим выражение:

\[
\texttt{if } x < 10 \texttt{ then } x + 1 \texttt{ else } 0.
\]

Пусть окружение типов:

\[
\Gamma = \{x : \texttt{Int}\}.
\]

Шаги проверки:

\begin{enumerate}
    \item Проверяется условие \(x < 10\).
    \item Переменная \(x\) имеет тип \(\texttt{Int}\).
    \item Константа \(10\) имеет тип \(\texttt{Int}\).
    \item Оператор \(<\) применён к двум значениям типа \(\texttt{Int}\), значит
    условие имеет тип \(\texttt{Bool}\).
    \item Проверяется ветвь \texttt{then}: \(x + 1\).
    \item \(x : \texttt{Int}\), \(1 : \texttt{Int}\), значит \(x + 1 : \texttt{Int}\).
    \item Проверяется ветвь \texttt{else}: \(0 : \texttt{Int}\).
    \item Обе ветви имеют тип \(\texttt{Int}\), значит всё выражение имеет тип
    \(\texttt{Int}\).
\end{enumerate}

Итог:

\[
\Gamma \vdash
\texttt{if } x < 10 \texttt{ then } x + 1 \texttt{ else } 0
:
\texttt{Int}.
\]

\subsubsection{Пример ошибки типов}

Выражение:

\[
\texttt{if } x < 10 \texttt{ then } x + 1 \texttt{ else } \texttt{false}.
\]

Пусть снова:

\[
\Gamma = \{x : \texttt{Int}\}.
\]

Тогда:

\[
x + 1 : \texttt{Int},
\]

но

\[
\texttt{false} : \texttt{Bool}.
\]

Типы ветвей не совпадают, поэтому выражение отвергается статической проверкой типов.

\subsection{Формальные правила типизации}

Для простого языка можно записать правила типизации в виде правил вывода.

\[
\frac{}{\Gamma \vdash n : \texttt{Int}}
\]

Целочисленная константа имеет тип \(\texttt{Int}\).

\[
\frac{}{\Gamma \vdash \texttt{true} : \texttt{Bool}}
\qquad
\frac{}{\Gamma \vdash \texttt{false} : \texttt{Bool}}
\]

Булевы константы имеют тип \(\texttt{Bool}\).

\[
\frac{x : \tau \in \Gamma}{\Gamma \vdash x : \tau}
\]

Если переменная \(x\) имеет тип \(\tau\) в окружении \(\Gamma\), то выражение \(x\)
имеет тип \(\tau\).

\[
\frac{\Gamma \vdash e_1 : \texttt{Int}
\qquad
\Gamma \vdash e_2 : \texttt{Int}}
{\Gamma \vdash e_1 + e_2 : \texttt{Int}}
\]

Сложение корректно только для двух целочисленных выражений.

\[
\frac{\Gamma \vdash e_1 : \texttt{Int}
\qquad
\Gamma \vdash e_2 : \texttt{Int}}
{\Gamma \vdash e_1 < e_2 : \texttt{Bool}}
\]

Сравнение двух целых чисел возвращает булево значение.

\[
\frac{
\Gamma \vdash e_1 : \texttt{Bool}
\qquad
\Gamma \vdash e_2 : \tau
\qquad
\Gamma \vdash e_3 : \tau
}
{
\Gamma \vdash \texttt{if } e_1 \texttt{ then } e_2 \texttt{ else } e_3 : \tau
}
\]

Условие должно иметь тип \(\texttt{Bool}\), а обе ветви должны иметь один и тот же
тип \(\tau\).

\subsection{Безопасность типов}

Одним из важнейших свойств хорошей системы типов является \textbf{безопасность типов}.

Интуитивно безопасность типов означает следующее:

\begin{quote}
Если программа успешно прошла проверку типов, то при выполнении она не попадёт в
состояние, где необходимо выполнить операцию над значениями неподходящих типов.
\end{quote}

Классическая формулировка безопасности типов обычно разбивается на две теоремы:

\begin{itemize}
    \item теорему прогресса;
    \item теорему сохранения типа.
\end{itemize}

Рассмотрим их для простого языка арифметических и булевых выражений.

\subsubsection{Синтаксис языка}

Пусть выражения задаются грамматикой:

\[
e ::= n
\mid \texttt{true}
\mid \texttt{false}
\mid e_1 + e_2
\mid e_1 < e_2
\mid \texttt{if } e_1 \texttt{ then } e_2 \texttt{ else } e_3.
\]

Значения:

\[
v ::= n \mid \texttt{true} \mid \texttt{false}.
\]

Типы:

\[
\tau ::= \texttt{Int} \mid \texttt{Bool}.
\]

Отношение малого шага вычисления обозначим:

\[
e \to e'.
\]

Оно означает, что выражение \(e\) за один шаг вычисления переходит в выражение \(e'\).

\subsubsection{Правила вычисления}

Для сложения:

\[
\frac{e_1 \to e_1'}{e_1 + e_2 \to e_1' + e_2}
\]

\[
\frac{e_2 \to e_2'}{n + e_2 \to n + e_2'}
\]

\[
n_1 + n_2 \to n_3,
\]

где \(n_3\) --- обычная сумма \(n_1\) и \(n_2\).

Для сравнения:

\[
\frac{e_1 \to e_1'}{e_1 < e_2 \to e_1' < e_2}
\]

\[
\frac{e_2 \to e_2'}{n < e_2 \to n < e_2'}
\]

\[
n_1 < n_2 \to \texttt{true},
\quad \text{если } n_1 < n_2,
\]

\[
n_1 < n_2 \to \texttt{false},
\quad \text{если } n_1 \ge n_2.
\]

Для условного оператора:

\[
\frac{e_1 \to e_1'}
{\texttt{if } e_1 \texttt{ then } e_2 \texttt{ else } e_3
\to
\texttt{if } e_1' \texttt{ then } e_2 \texttt{ else } e_3}
\]

\[
\texttt{if true then } e_2 \texttt{ else } e_3 \to e_2,
\]

\[
\texttt{if false then } e_2 \texttt{ else } e_3 \to e_3.
\]

\subsubsection{Теорема прогресса}

\textbf{Теорема прогресса.}
Если выражение \(e\) замкнуто, то есть не содержит свободных переменных, и

\[
\vdash e : \tau,
\]

то либо \(e\) является значением, либо существует выражение \(e'\), такое что

\[
e \to e'.
\]

Иначе говоря, хорошо типизированное замкнутое выражение либо уже является
результатом вычисления, либо может сделать следующий шаг вычисления.

\textbf{Доказательство.}

Доказательство проводится структурной индукцией по выводу типизации
\(\vdash e : \tau\).

\textbf{Случай 1.} \(e = n\).

По определению \(n\) является значением. Следовательно, утверждение выполнено.

\textbf{Случай 2.} \(e = \texttt{true}\) или \(e = \texttt{false}\).

Булевы константы являются значениями. Следовательно, утверждение выполнено.

\textbf{Случай 3.} \(e = e_1 + e_2\).

По правилу типизации сложения:

\[
\vdash e_1 : \texttt{Int},
\qquad
\vdash e_2 : \texttt{Int}.
\]

По индукционному предположению для \(e_1\) имеем:

\[
e_1 \text{ является значением}
\quad \text{или} \quad
\exists e_1' : e_1 \to e_1'.
\]

Если \(e_1 \to e_1'\), то по правилу вычисления:

\[
e_1 + e_2 \to e_1' + e_2.
\]

Значит, выражение делает шаг.

Если \(e_1\) является значением, то из \(\vdash e_1 : \texttt{Int}\) следует, что
\(e_1\) имеет вид некоторого числа \(n_1\). Теперь применим индукционное предположение
к \(e_2\).

Если \(e_2 \to e_2'\), то:

\[
n_1 + e_2 \to n_1 + e_2'.
\]

Если \(e_2\) является значением, то из \(\vdash e_2 : \texttt{Int}\) следует, что
\(e_2 = n_2\). Тогда:

\[
n_1 + n_2 \to n_3,
\]

где \(n_3\) --- сумма чисел \(n_1\) и \(n_2\). Следовательно, выражение делает шаг.

\textbf{Случай 4.} \(e = e_1 < e_2\).

Рассуждение аналогично случаю сложения. По правилу типизации:

\[
\vdash e_1 : \texttt{Int},
\qquad
\vdash e_2 : \texttt{Int}.
\]

Если \(e_1\) делает шаг, то всё выражение делает шаг. Если \(e_1\) является числом
\(n_1\), рассматривается \(e_2\). Если \(e_2\) делает шаг, то всё выражение делает шаг.
Если \(e_2\) является числом \(n_2\), то выражение \(n_1 < n_2\) вычисляется либо в
\(\texttt{true}\), либо в \(\texttt{false}\). Следовательно, шаг существует.

\textbf{Случай 5.}
\[
e = \texttt{if } e_1 \texttt{ then } e_2 \texttt{ else } e_3.
\]

По правилу типизации условного выражения:

\[
\vdash e_1 : \texttt{Bool},
\qquad
\vdash e_2 : \tau,
\qquad
\vdash e_3 : \tau.
\]

По индукционному предположению для \(e_1\):

\[
e_1 \text{ является значением}
\quad \text{или} \quad
\exists e_1' : e_1 \to e_1'.
\]

Если \(e_1 \to e_1'\), то:

\[
\texttt{if } e_1 \texttt{ then } e_2 \texttt{ else } e_3
\to
\texttt{if } e_1' \texttt{ then } e_2 \texttt{ else } e_3.
\]

Если \(e_1\) является значением, то из \(\vdash e_1 : \texttt{Bool}\) следует, что
\(e_1\) равно либо \(\texttt{true}\), либо \(\texttt{false}\).

Если \(e_1 = \texttt{true}\), то:

\[
\texttt{if true then } e_2 \texttt{ else } e_3 \to e_2.
\]

Если \(e_1 = \texttt{false}\), то:

\[
\texttt{if false then } e_2 \texttt{ else } e_3 \to e_3.
\]

Следовательно, выражение либо является значением, либо может сделать шаг.
Теорема доказана.

\subsubsection{Теорема сохранения типа}

\textbf{Теорема сохранения типа.}
Если

\[
\vdash e : \tau
\]

и

\[
e \to e',
\]

то

\[
\vdash e' : \tau.
\]

Иначе говоря, один шаг вычисления не меняет тип выражения.

\textbf{Доказательство.}

Доказательство проводится индукцией по выводу отношения вычисления \(e \to e'\).

\textbf{Случай 1.}
\[
e = e_1 + e_2,
\qquad
e_1 \to e_1',
\qquad
e' = e_1' + e_2.
\]

Из типизации \(e_1 + e_2\) следует:

\[
\vdash e_1 : \texttt{Int},
\qquad
\vdash e_2 : \texttt{Int},
\qquad
\tau = \texttt{Int}.
\]

По индукционному предположению из \(e_1 \to e_1'\) и
\(\vdash e_1 : \texttt{Int}\) получаем:

\[
\vdash e_1' : \texttt{Int}.
\]

Тогда по правилу типизации сложения:

\[
\vdash e_1' + e_2 : \texttt{Int}.
\]

Следовательно:

\[
\vdash e' : \tau.
\]

\textbf{Случай 2.}
\[
e = n + e_2,
\qquad
e_2 \to e_2',
\qquad
e' = n + e_2'.
\]

Из типизации \(n + e_2\) следует:

\[
\vdash n : \texttt{Int},
\qquad
\vdash e_2 : \texttt{Int}.
\]

По индукционному предположению:

\[
\vdash e_2' : \texttt{Int}.
\]

Следовательно:

\[
\vdash n + e_2' : \texttt{Int}.
\]

\textbf{Случай 3.}
\[
e = n_1 + n_2,
\qquad
e' = n_3,
\]

где \(n_3\) --- сумма \(n_1\) и \(n_2\).

Из правила типизации констант:

\[
\vdash n_3 : \texttt{Int}.
\]

Так как исходное выражение \(n_1 + n_2\) тоже имеет тип \(\texttt{Int}\), тип
сохраняется.

\textbf{Случай 4.}
\[
e = e_1 < e_2.
\]

Рассуждение аналогично сложению. Если вычисляется левый или правый операнд, то по
индукционному предположению его тип \(\texttt{Int}\) сохраняется. Если оба операнда
являются числами, то результатом вычисления является \(\texttt{true}\) или
\(\texttt{false}\), то есть выражение типа \(\texttt{Bool}\). Следовательно, тип
сохраняется.

\textbf{Случай 5.}
\[
e =
\texttt{if } e_1 \texttt{ then } e_2 \texttt{ else } e_3,
\qquad
e_1 \to e_1'.
\]

Тогда:

\[
e' =
\texttt{if } e_1' \texttt{ then } e_2 \texttt{ else } e_3.
\]

Из типизации исходного выражения следует:

\[
\vdash e_1 : \texttt{Bool},
\qquad
\vdash e_2 : \tau,
\qquad
\vdash e_3 : \tau.
\]

По индукционному предположению:

\[
\vdash e_1' : \texttt{Bool}.
\]

Следовательно, по правилу типизации условного выражения:

\[
\vdash
\texttt{if } e_1' \texttt{ then } e_2 \texttt{ else } e_3
:
\tau.
\]

\textbf{Случай 6.}
\[
e = \texttt{if true then } e_2 \texttt{ else } e_3,
\qquad
e' = e_2.
\]

Из типизации исходного выражения следует:

\[
\vdash e_2 : \tau,
\qquad
\vdash e_3 : \tau.
\]

Значит:

\[
\vdash e' : \tau.
\]

\textbf{Случай 7.}
\[
e = \texttt{if false then } e_2 \texttt{ else } e_3,
\qquad
e' = e_3.
\]

Из типизации исходного выражения следует:

\[
\vdash e_2 : \tau,
\qquad
\vdash e_3 : \tau.
\]

Значит:

\[
\vdash e' : \tau.
\]

Все возможные правила вычисления рассмотрены. Теорема доказана.

\subsubsection{Следствие: безопасность типов}

\textbf{Следствие.}
Если замкнутое выражение \(e\) хорошо типизировано, то его вычисление либо
завершается значением того же типа, либо может продолжаться бесконечно, но не
застревает из-за ошибки типов.

\textbf{Обоснование.}

Из теоремы сохранения типа следует, что при каждом шаге вычисления тип выражения
сохраняется. Из теоремы прогресса следует, что хорошо типизированное выражение,
которое не является значением, может сделать следующий шаг. Поэтому невозможно
получить состояние, в котором выражение ещё не является значением, но дальнейшее
вычисление невозможно из-за некорректной операции над типами.

\subsection{Статическая типизация и вывод типов}

Статическая типизация не обязательно требует, чтобы программист явно писал тип
каждой переменной. Компилятор может использовать \textbf{вывод типов}.

\textbf{Вывод типов} --- это автоматическое определение типов выражений на основе
контекста их использования.

Например, в OCaml:

\begin{verbatim}
let f x = x + 1
\end{verbatim}

Компилятор выводит:

\[
f : \texttt{int} \to \texttt{int}.
\]

Поскольку операция \texttt{+} применима к целым числам, аргумент \(x\) должен иметь
тип \texttt{int}, а результат также имеет тип \texttt{int}.

\subsubsection{Идея алгоритма унификации}

Во многих языках с выводом типов используется \textbf{унификация}.

\textbf{Унификация} --- это процесс нахождения подстановки для типовых переменных,
при которой два типа становятся одинаковыми.

Например, требуется унифицировать:

\[
\alpha \to \texttt{Int}
\]

и

\[
\texttt{Bool} \to \beta.
\]

Тогда получаем ограничения:

\[
\alpha = \texttt{Bool},
\qquad
\beta = \texttt{Int}.
\]

Искомая подстановка:

\[
\{\alpha \mapsto \texttt{Bool}, \beta \mapsto \texttt{Int}\}.
\]

\subsubsection{Минимальный алгоритм унификации типов}

Пусть типы имеют вид:

\[
\tau ::= \texttt{Int}
\mid \texttt{Bool}
\mid \alpha
\mid \tau_1 \to \tau_2.
\]

Алгоритм унификации множества уравнений типов:

\[
\{\tau_1 = \sigma_1, \ldots, \tau_n = \sigma_n\}
\]

работает следующим образом:

\begin{enumerate}
    \item Если множество уравнений пусто, вернуть найденную подстановку.
    \item Взять одно уравнение \(\tau = \sigma\).
    \item Если \(\tau\) и \(\sigma\) одинаковые атомарные типы, отбросить уравнение.
    \item Если \(\tau\) является типовой переменной \(\alpha\), заменить \(\alpha\)
    на \(\sigma\) во всех оставшихся уравнениях и добавить подстановку
    \(\alpha \mapsto \sigma\).
    \item Если \(\sigma\) является типовой переменной \(\alpha\), действовать
    аналогично.
    \item Если
    \[
    \tau = \tau_1 \to \tau_2,
    \qquad
    \sigma = \sigma_1 \to \sigma_2,
    \]
    заменить уравнение \(\tau = \sigma\) двумя уравнениями:
    \[
    \tau_1 = \sigma_1,
    \qquad
    \tau_2 = \sigma_2.
    \]
    \item Если ни одно правило не применимо, унификация невозможна: имеется ошибка
    типов.
\end{enumerate}

Для корректного алгоритма также необходима \textbf{проверка вхождения}: нельзя
унифицировать типовую переменную \(\alpha\) с типом, который уже содержит \(\alpha\),
например:

\[
\alpha = \alpha \to \texttt{Int}.
\]

Иначе возник бы бесконечный тип.

\subsubsection{Пример унификации}

Пусть нужно унифицировать:

\[
\alpha \to \alpha
\]

и

\[
\texttt{Int} \to \beta.
\]

Шаги:

\begin{enumerate}
    \item Исходное уравнение:
    \[
    \alpha \to \alpha = \texttt{Int} \to \beta.
    \]
    \item Разбиваем его на два уравнения:
    \[
    \alpha = \texttt{Int},
    \qquad
    \alpha = \beta.
    \]
    \item Из первого уравнения получаем:
    \[
    \alpha \mapsto \texttt{Int}.
    \]
    \item Подставляем это во второе уравнение:
    \[
    \texttt{Int} = \beta.
    \]
    \item Получаем:
    \[
    \beta \mapsto \texttt{Int}.
    \]
\end{enumerate}

Итоговая подстановка:

\[
\{\alpha \mapsto \texttt{Int}, \beta \mapsto \texttt{Int}\}.
\]

\subsection{Преимущества и недостатки статической типизации}

\subsubsection{Преимущества}

\begin{itemize}
    \item \textbf{Раннее обнаружение ошибок.}
    Ошибки типов находятся до запуска программы.

    \item \textbf{Документирование интерфейсов.}
    Типы функций и методов описывают, какие данные ожидаются и какие возвращаются.

    \item \textbf{Улучшенная поддержка IDE.}
    Статические типы помогают автодополнению, рефакторингу, навигации по коду.

    \item \textbf{Оптимизация.}
    Компилятор может использовать информацию о типах для генерации более эффективного
    кода.

    \item \textbf{Надёжность крупных систем.}
    При изменении интерфейса компилятор помогает найти все места, которые нужно
    исправить.
\end{itemize}

\subsubsection{Недостатки}

\begin{itemize}
    \item \textbf{Большая формальность.}
    Иногда требуется писать больше аннотаций типов.

    \item \textbf{Меньшая гибкость.}
    Некоторые динамические приёмы программирования сложнее выразить.

    \item \textbf{Сложность системы типов.}
    В современных языках типы могут быть весьма сложными: обобщения, зависимые типы,
    высшие типы, трейты, ограничения и т.д.

    \item \textbf{Время компиляции.}
    Большие проекты могут компилироваться долго.
\end{itemize}

\subsection{Преимущества и недостатки динамической типизации}

\subsubsection{Преимущества}

\begin{itemize}
    \item \textbf{Быстрая разработка.}
    Можно писать меньше служебного кода.

    \item \textbf{Гибкость.}
    Удобно создавать прототипы, скрипты, обработчики данных.

    \item \textbf{Интерактивность.}
    Многие динамические языки имеют REPL-среды, где можно быстро выполнять фрагменты
    кода.

    \item \textbf{Простота метапрограммирования.}
    Во многих динамических языках легко изменять объекты, классы и функции во время
    выполнения.
\end{itemize}

\subsubsection{Недостатки}

\begin{itemize}
    \item \textbf{Ошибки типов могут обнаружиться поздно.}
    Если участок кода редко выполняется, ошибка может долго оставаться незамеченной.

    \item \textbf{Сложность поддержки больших проектов.}
    Без дополнительных соглашений и тестов труднее отслеживать контракты между
    модулями.

    \item \textbf{Меньше возможностей для оптимизации.}
    Реализация языка должна учитывать, что тип значения может быть известен только
    во время выполнения.

    \item \textbf{Необходимость большего числа тестов.}
    Тесты часто компенсируют отсутствие статической проверки.
\end{itemize}

\subsection{Постепенная типизация}

\textbf{Постепенная типизация} --- это подход, при котором язык допускает как
типизированные, так и нетипизированные части программы.

Примером является TypeScript по отношению к JavaScript.

JavaScript:

\begin{verbatim}
function add(x, y) {
    return x + y;
}
\end{verbatim}

TypeScript:

\begin{verbatim}
function add(x: number, y: number): number {
    return x + y;
}
\end{verbatim}

Преимущество постепенной типизации состоит в том, что большой динамический проект
можно переводить на статические типы постепенно, не переписывая всё сразу.

Также элементы постепенной типизации есть в современных версиях Python через
аннотации типов:

\begin{verbatim}
def add(x: int, y: int) -> int:
    return x + y
\end{verbatim}

Сами по себе аннотации типов в Python не меняют семантику выполнения программы,
но могут проверяться внешними инструментами, например статическими анализаторами.

\subsection{Языки сценариев}

\subsubsection{Определение}

\textbf{Язык сценариев}, или \textbf{скриптовый язык}, --- это язык программирования,
ориентированный на автоматизацию задач, управление другими программами, связывание
компонентов, быструю разработку и выполнение небольших или средних программ без
сложного этапа компиляции.

Сценарий, или скрипт, обычно представляет собой программу, которая:

\begin{itemize}
    \item управляет уже существующими компонентами;
    \item автоматизирует последовательность действий;
    \item обрабатывает текстовые данные;
    \item запускает внешние команды;
    \item конфигурирует приложение;
    \item расширяет поведение некоторой системы.
\end{itemize}

Примеры скриптовых языков:

\begin{itemize}
    \item Unix shell;
    \item Bash;
    \item Python;
    \item Perl;
    \item Ruby;
    \item Tcl;
    \item Lua;
    \item JavaScript;
    \item PowerShell;
    \item PHP.
\end{itemize}

\subsubsection{Характерные свойства языков сценариев}

Для языков сценариев часто характерны следующие свойства:

\begin{itemize}
    \item \textbf{Интерпретируемое выполнение.}
    Программа часто выполняется интерпретатором без отдельной явной компиляции.

    \item \textbf{Динамическая типизация.}
    Многие скриптовые языки являются динамически типизированными.

    \item \textbf{Высокоуровневые встроенные структуры данных.}
    Например, списки, словари, множества, строки, регулярные выражения.

    \item \textbf{Удобная работа с текстом.}
    Исторически многие скриптовые языки применялись для обработки файлов, журналов,
    отчётов и текстовых потоков.

    \item \textbf{Простое взаимодействие с ОС.}
    Удобный запуск команд, работа с файлами, переменными окружения, процессами.

    \item \textbf{Краткость программ.}
    Типичные задачи автоматизации можно записывать небольшим количеством строк.

    \item \textbf{Использование в качестве ``клея''.}
    Скриптовый язык часто связывает компоненты, написанные на других языках.
\end{itemize}

\subsubsection{Скриптовые языки и системные языки}

Условно можно противопоставить:

\begin{itemize}
    \item \textbf{системные языки};
    \item \textbf{скриптовые языки}.
\end{itemize}

\textbf{Системные языки} ориентированы на эффективную реализацию низкоуровневых
механизмов, библиотек, операционных систем, компиляторов, баз данных, виртуальных
машин и производительных приложений.

Примеры:

\begin{itemize}
    \item C;
    \item C++;
    \item Rust;
    \item Zig;
    \item Go в некоторых областях.
\end{itemize}

\textbf{Скриптовые языки} ориентированы на быстрое описание логики, автоматизацию и
интеграцию готовых компонентов.

Примеры:

\begin{itemize}
    \item Python;
    \item Bash;
    \item Perl;
    \item Ruby;
    \item Lua;
    \item Tcl;
    \item JavaScript.
\end{itemize}

Например, вычислительно сложная библиотека машинного обучения может быть написана
на C++ или CUDA, а пользовательский код, который вызывает эту библиотеку, --- на
Python.

\subsubsection{Пример скрипта на Bash}

\begin{verbatim}
#!/bin/bash

for file in *.txt
do
    echo "Processing $file"
    wc -l "$file"
done
\end{verbatim}

Этот скрипт перебирает все файлы с расширением \texttt{.txt} в текущем каталоге и
выводит число строк в каждом файле.

Шаги выполнения:

\begin{enumerate}
    \item Оболочка находит все имена файлов, подходящие под шаблон \texttt{*.txt}.
    \item Переменная \texttt{file} последовательно принимает каждое найденное имя.
    \item Команда \texttt{echo} выводит сообщение.
    \item Команда \texttt{wc -l} подсчитывает количество строк в файле.
\end{enumerate}

\subsubsection{Пример скрипта на Python}

\begin{verbatim}
import os

for name in os.listdir("."):
    if name.endswith(".txt"):
        print(name)
\end{verbatim}

Скрипт выводит имена всех файлов с расширением \texttt{.txt} в текущем каталоге.

Шаги выполнения:

\begin{enumerate}
    \item Импортируется модуль \texttt{os}.
    \item Функция \texttt{os.listdir(".")} возвращает список имён файлов и каталогов.
    \item Цикл перебирает эти имена.
    \item Метод \texttt{endswith(".txt")} проверяет расширение.
    \item Если условие истинно, имя выводится на экран.
\end{enumerate}

\subsubsection{Пример использования скриптового языка как ``клея''}

Пусть есть внешняя программа \texttt{convert}, преобразующая изображения. Python-скрипт
может автоматически вызвать её для набора файлов:

\begin{verbatim}
import subprocess
from pathlib import Path

for path in Path(".").glob("*.png"):
    output = path.with_suffix(".jpg")
    subprocess.run(["convert", str(path), str(output)])
\end{verbatim}

Здесь Python не обязательно сам реализует алгоритм преобразования изображения.
Он управляет внешней программой и автоматизирует пакетную обработку файлов.

\subsection{Интерпретация и компиляция}

Скриптовые языки часто называют интерпретируемыми, но это упрощение.

\textbf{Интерпретация} --- выполнение программы непосредственно интерпретатором.

\textbf{Компиляция} --- перевод программы из одного языка в другой, обычно из
высокоуровневого языка в машинный код, байт-код или промежуточное представление.

Многие современные реализации используют смешанные подходы:

\begin{itemize}
    \item Python компилирует исходный код в байт-код, который затем исполняется
    виртуальной машиной;
    \item Java компилируется в байт-код JVM;
    \item JavaScript в современных движках часто использует JIT-компиляцию;
    \item C и C++ обычно компилируются заранее в машинный код;
    \item Java и C\# сочетают байт-код и JIT-компиляцию.
\end{itemize}

Следовательно, нельзя строго отождествлять:

\[
\text{скриптовый язык} \neq \text{обязательно интерпретируемый язык}.
\]

И также:

\[
\text{компилируемый язык} \neq \text{обязательно статически типизированный язык}.
\]

Например, Java является статически типизированной и компилируется в байт-код, а
JavaScript является динамически типизированным, но современные реализации активно
используют компиляцию во время выполнения.

\subsection{Типизация и выполнение: независимые классификации}

Языки можно классифицировать по нескольким независимым признакам:

\begin{enumerate}
    \item Когда проверяются типы:
    \begin{itemize}
        \item статически;
        \item динамически.
    \end{itemize}

    \item Насколько строго язык обращается с типами:
    \begin{itemize}
        \item сильная типизация;
        \item слабая типизация.
    \end{itemize}

    \item Нужно ли явно писать типы:
    \begin{itemize}
        \item явная типизация;
        \item неявная типизация;
        \item частичный вывод типов.
    \end{itemize}

    \item Как выполняется программа:
    \begin{itemize}
        \item интерпретация;
        \item предварительная компиляция;
        \item компиляция в байт-код;
        \item JIT-компиляция.
    \end{itemize}

    \item Какова основная область применения:
    \begin{itemize}
        \item системное программирование;
        \item прикладное программирование;
        \item сценарии и автоматизация;
        \item научные вычисления;
        \item веб-разработка;
        \item встроенные системы.
    \end{itemize}
\end{enumerate}

Эти признаки не сводятся друг к другу. Например:

\begin{itemize}
    \item Python --- динамически типизированный, обычно сильно типизированный,
    часто используется как скриптовый язык;
    \item Java --- статически типизированный, сильно типизированный, компилируется
    в байт-код;
    \item C --- статически типизированный, но допускает низкоуровневые преобразования,
    используется как системный язык;
    \item JavaScript --- динамически типизированный, имеет слабые аспекты из-за
    неявных преобразований, широко используется как скриптовый язык;
    \item Haskell --- статически типизированный, сильно типизированный, с развитым
    выводом типов.
\end{itemize}

\subsection{Минимальные примеры классификации языков}

\subsubsection{Python}

\begin{verbatim}
x = 1
x = "one"
\end{verbatim}

Переменная \texttt{x} может ссылаться на значения разных типов, следовательно,
типизация динамическая.

\begin{verbatim}
"1" + 2
\end{verbatim}

Такая операция вызывает ошибку, следовательно, язык не выполняет произвольного
неявного преобразования. Поэтому Python обычно относят к сильно типизированным
динамическим языкам.

\subsubsection{Java}

\begin{verbatim}
int x = 1;
x = "one";
\end{verbatim}

Ошибка обнаруживается компилятором. Java является статически типизированным языком.

\subsubsection{JavaScript}

\begin{verbatim}
let x = "1" + 2;  // "12"
let y = "1" - 2;  // -1
\end{verbatim}

Типы проверяются во время выполнения, а язык выполняет неявные преобразования.
Поэтому JavaScript является динамически типизированным языком со слабыми аспектами
типизации.

\subsubsection{C}

\begin{verbatim}
int x = 10;
double y = x;
\end{verbatim}

Компилятор знает типы заранее, следовательно, язык статически типизирован.

Но при этом C допускает низкоуровневые преобразования:

\begin{verbatim}
int x = 10;
char *p = (char *)&x;
\end{verbatim}

Такие конструкции ослабляют безопасность типов.

\subsection{Итоговая схема}

\begin{center}
\begin{tabular}{|p{0.26\textwidth}|p{0.34\textwidth}|p{0.32\textwidth}|}
\hline
\textbf{Понятие} & \textbf{Смысл} & \textbf{Пример} \\
\hline
Статическая типизация &
Типы проверяются до выполнения программы &
Java, C, Rust, Haskell \\
\hline
Динамическая типизация &
Типы проверяются во время выполнения &
Python, Ruby, JavaScript \\
\hline
Сильная типизация &
Язык запрещает опасные или бессмысленные операции над несовместимыми типами &
Python, Java, Haskell \\
\hline
Слабая типизация &
Язык допускает широкие неявные преобразования или низкоуровневые интерпретации данных &
JavaScript, C в некоторых аспектах \\
\hline
Скриптовый язык &
Язык для автоматизации, управления приложениями, связывания компонентов и быстрой разработки &
Bash, Python, Perl, Lua, Tcl \\
\hline
\end{tabular}
\end{center}

\subsection{Краткие выводы}

\begin{enumerate}
    \item Статическая типизация проверяет типы до выполнения программы, динамическая ---
    во время выполнения.

    \item Статическая типизация помогает раньше находить ошибки, улучшает поддержку
    больших проектов и даёт больше возможностей для оптимизации.

    \item Динамическая типизация удобна для быстрой разработки, прототипирования,
    интерактивной работы и метапрограммирования.

    \item Сильная и слабая типизация описывают не момент проверки типов, а строгость
    правил обращения с типами.

    \item Сильная типизация не равна статической: Python динамический, но обычно
    считается сильно типизированным.

    \item Слабая типизация не равна динамической: C статически типизирован, но допускает
    низкоуровневые преобразования, ослабляющие безопасность типов.

    \item Скриптовые языки обычно ориентированы на автоматизацию, обработку данных,
    управление программами и связывание компонентов.

    \item Скриптовость, интерпретируемость, динамическая типизация и слабая типизация
    --- разные свойства, которые часто встречаются вместе, но логически независимы.
\end{enumerate}

\end{document}
